\section{Emergent Kinematics on the $\Tdiamond$ Lattice}
\label{sec:kinematics}

Momentum, inertia, and force are not postulated in the $\unity$ framework.
They emerge from the geometry of phase interference across the bipartite
sublattice structure.  This section derives the kinematic relations
from the Dirac tick rule and connects each result to the corresponding
numerical verification in the experiment suite.

% ── Subsection 1 ─────────────────────────────────────────────────────────────
\subsection{The Bipartite Velocity}
\label{subsec:bipartite_velocity}

A massive Dirac fermion~\cite{dirac1928} on $\Tdiamond$ carries two spinor components:
$\psi_R$ on the RGB sublattice (even ticks) and $\psi_L$ on the CMY
sublattice (odd ticks).  The $\unity$ constraint requires
\begin{equation}
    \sum_{\mathbf{x}} \bigl(|\psi_R(\mathbf{x})|^2 +
                            |\psi_L(\mathbf{x})|^2\bigr) = 1.
    \label{eq:unity_spinor}
\end{equation}
Over one complete two-tick cycle the kinetic fraction
$p_\mathrm{hop} = \cos^2(\omega/2)$ is distributed across the six
basis vectors, split equally between the three RGB directions
($\mathbf{V}_{1\text{--}3}$, even tick) and the three CMY directions
($-\mathbf{V}_{1\text{--}3}$, odd tick):
\begin{equation}
    1 = p_\mathrm{stay} + P_\mathrm{fwd} + P_\mathrm{bwd},
    \qquad
    p_\mathrm{stay} = \sin^2\!\tfrac{\omega}{2},
    \quad
    P_\mathrm{fwd} + P_\mathrm{bwd} = \cos^2\!\tfrac{\omega}{2}.
    \label{eq:probability_partition}
\end{equation}
Here $P_\mathrm{fwd} = \sum_{i=1}^{3} p(\mathbf{V}_i)$ is the total
probability flowing along the RGB sublattice and
$P_\mathrm{bwd} = \sum_{i=1}^{3} p(-\mathbf{V}_i)$ along CMY.
The lattice speed limit $c = 1$ node/tick gives the net velocity
\begin{equation}
    v = c\,(P_\mathrm{fwd} - P_\mathrm{bwd}).
    \label{eq:velocity}
\end{equation}
For a packet at rest the two sublattices are populated symmetrically,
$P_\mathrm{fwd} = P_\mathrm{bwd} = \tfrac{1}{2}\cos^2(\omega/2)$,
so $v = 0$ and the amplitude oscillates between $\psi_R$ and $\psi_L$
without net displacement — the lattice realisation of Zitterbewegung~\cite{schrodinger1930}.

% ── Subsection 2 ─────────────────────────────────────────────────────────────
\subsection{Phase Gradient as Momentum}
\label{subsec:momentum_emergence}

A uniform spatial phase gradient with advance $\theta$ per lattice edge
breaks the RGB/CMY symmetry.  In the implementation this gradient
appears as the per-direction phase advance
\begin{equation}
    \delta_p(\mathbf{V}_i, \mathbf{x})
        = \arg\psi(\mathbf{x} + \mathbf{V}_i) - \arg\psi(\mathbf{x}),
    \label{eq:delta_p}
\end{equation}
which is the quantity \texttt{delta\_p} computed in
\texttt{CausalSession.\_kinetic\_hop}.  Only positive advances
contribute to the hop weight, so amplitude flows preferentially in the
direction of increasing phase — the lattice realisation of the
de~Broglie relation~\cite{debroglie1924}.

To derive $v(\theta)$ analytically, consider a packet with mass
parameter $\omega$ (the instruction frequency) subject to a gradient
$\theta$ directed along the net RGB axis.  On the even (RGB) tick
$\psi_R$ hops forward: the phase gradient $+\theta$ reduces the
effective mismatch from $\omega$ to $\omega - \theta$, increasing the
kinetic weight.  On the odd (CMY) tick $\psi_L$ hops backward: the
reversed gradient adds $+\theta$ to the mismatch, decreasing the
kinetic weight:
\begin{align}
    P_\mathrm{fwd} &= \tfrac{1}{2}\cos^2\!\tfrac{\omega - \theta}{2},
    \qquad (\psi_R\text{ hop, even tick})
    \label{eq:Pfwd} \\
    P_\mathrm{bwd} &= \tfrac{1}{2}\cos^2\!\tfrac{\omega + \theta}{2}.
    \qquad (\psi_L\text{ hop, odd tick})
    \label{eq:Pbwd}
\end{align}
Substituting into Eq.~\eqref{eq:velocity} and applying
$\cos^2\alpha - \cos^2\beta = -\sin(\alpha+\beta)\sin(\alpha-\beta)$
yields
\begin{equation}
    v = \tfrac{c}{2}\,\sin(\omega)\,\sin(\theta).
    \label{eq:v_omega_theta}
\end{equation}
Table~\ref{tab:mass_velocity} illustrates how $p_\mathrm{stay}$,
$m_\mathrm{inertial}$, and the resulting velocity vary across
representative values of $\omega$:
\begin{table}[h]
\centering
\begin{tabular}{cccc}
\hline
$\omega$ & $p_\mathrm{stay} = \sin^2(\omega/2)$ &
           $m_\mathrm{inertial} = \sin(\omega)/2$ &
           $v/c$ at $\theta = 0.1$ \\
\hline
$0$       & $0$     & $0$     & $1$ \quad (photon limit) \\
$\pi/4$   & $0.146$ & $0.354$ & $0.035$ \\
$\pi/2$   & $0.500$ & $0.500$ & $0.050$ \\
$\pi$     & $1.000$ & $0$     & $0$ \quad (frozen, max mass) \\
\hline
\end{tabular}
\caption{Selected values of the two mass-related quantities and the
  resulting lattice velocity at a small phase gradient $\theta = 0.1$.
  The photon row ($\omega = 0$) is handled by the massless limit
  (Section~\ref{sec:causal_cone}); the other rows follow
  Eq.~\eqref{eq:v_omega_theta}.}
\label{tab:mass_velocity}
\end{table}
In the macroscopic limit $\theta \ll 1$ (gradient small compared to
$2\pi$), $\sin\theta \approx \theta$.  Identifying the lattice momentum
$p \equiv m\,v / c$ and the effective mass $m \equiv \sin(\omega)/2$\footnote{%
  The inertial mass $m_\mathrm{inertial} = \sin(\omega)/2$ is \emph{not} the
  same as the residence probability $p_\mathrm{stay} = \sin^2(\omega/2)$.
  They coincide only at $\omega = \pi/2$, the unique frequency at which
  the particle spends exactly half its time at rest.  At $\omega = \pi$
  (maximum residence) the particle is frozen, $m_\mathrm{inertial} = 0$,
  while $p_\mathrm{stay} = 1$.}
recovers the classical kinematic relation
\begin{equation}
    p = m\,\theta \cdot c,
    \label{eq:momentum}
\end{equation}
where $\theta$ is the phase advance per lattice edge — the discrete
analogue of the wavenumber $k$.  Momentum is thus strictly the
cross-term between the particle's internal mass-phase $\omega$ and the
spatial phase gradient of the bipartite vacuum.

\paragraph{Massless limit.}
Equation~\eqref{eq:v_omega_theta} gives $v = 0$ when $\omega = 0$,
which is not the photon behaviour.  The massless case is qualitatively
distinct: with $\omega = 0$ the residence probability vanishes,
$p_\mathrm{stay} = 0$, and the amplitude hops entirely via the
directional weighting of Eq.~\eqref{eq:delta_p}.  The hop is strictly
one-way along the gradient ($\delta_p > 0$ only), so $P_\mathrm{bwd}
= 0$ and $v = c$ regardless of $\theta$.  The two-sublattice
derivation above applies to massive particles ($\omega > 0$); the
photon is the $\omega \to 0$ limit treated separately in
Section~\ref{sec:causal_cone}.

\paragraph{Numerical verification.}
Experiment~\texttt{exp\_01\_inertia.py} initialises a Gaussian wavepacket
($\sigma = 3$, $\omega = 0.2$) with momentum $\mathbf{k} = (k,k,k)$,
$k = 0.3/\sqrt{3}$, corresponding to a phase advance
$\theta = \mathbf{k} \cdot \mathbf{V}_1 = 0.3\sqrt{3} \approx 0.520$
rad per hop along the $\mathbf{V}_1 = (1,1,1)$ direction.
Equation~\eqref{eq:v_omega_theta} predicts
\[
    v_x = \tfrac{1}{2}\sin(0.2)\sin(0.3\sqrt{3})
        \approx 0.0493 \text{ nodes/tick}.
\]
The measured x-component fit over 12 ticks gives
$v_x = 0.0464$ nodes/tick — a 6\% discrepancy attributable to
finite packet width and the short run length.  The heavy particle
($\omega = 0.5$, same $\mathbf{k}$) drifts measurably slower,
confirming that $\sin(\omega)$ governs inertial resistance: a larger
phase cost per hop reduces the net kinetic probability imbalance.

% ── Subsection 3 ─────────────────────────────────────────────────────────────
\subsection{Acceleration and Force}
\label{subsec:force}

Acceleration requires a time-varying gradient.  Taking the discrete
time derivative of Eq.~\eqref{eq:v_omega_theta} with $\omega$ held
constant:
\begin{equation}
    a = \frac{\Delta v}{\Delta t}
      = \frac{c}{2}\,\sin(\omega)\,\cos(\theta)\,\frac{\Delta\theta}{\Delta t}.
    \label{eq:acceleration}
\end{equation}
For small gradients ($\cos\theta \approx 1$) this reduces to the
lattice analogue of Newton's second law~\cite{newton1687}:
\begin{equation}
    a = \frac{c}{2m}\,F_\theta,
    \qquad
    F_\theta \equiv \frac{\Delta\theta}{\Delta t}.
    \label{eq:newton}
\end{equation}
The quantity $F_\theta$ — the rate at which the spatial phase gradient
changes per tick — is the lattice definition of force.  A gravitational
or electromagnetic field acts by twisting the phase of the vacuum over
time, unbalancing $P_\mathrm{fwd}$ and $P_\mathrm{bwd}$ and
continuously converting $p_\mathrm{stay}$ into directed kinetic
probability.  The mechanism by which the Coulomb potential produces
$F_\theta$ in the hydrogen orbital is analysed in
Section~\ref{sec:hydrogen}.
